\section{Evaluation Protocol}
\subsection{Separation of capture and assessment}
The recorder receives measurements and capture scores, never event labels. Packet
bytes are written, decoded, and only then evaluated. In the synthetic check,
capture uses the largest absolute first difference; the receiver uses absolute
amplitude. This avoids literally reusing the same score, but both tests are simple
and do not establish independence from shared modeling assumptions. The mission
study must include a separately fitted forecasting or relationship-based receiver,
and ideally a blinded engineering assessment.

For event interval $A_e$, raw coverage is the fraction of available source
observations inside $A_e$ retained exactly. Point events use their exact native
instant. Pre-event coverage uses available observations in a stated interval before
the annotated start. Raw onset error is the delay to the first retained event
sample, when one exists. We separately record whether a receiver finds evidence
and when the relevant packet becomes available.

A summary can overlap an event only at bucket resolution. A large extremum inside
an overlapping bucket is therefore \emph{coarse evidence}, not proof that its
source time matches the annotation. Summary evidence must not be presented as
precise onset recovery or raw-sample coverage. Events without available raw
observations require explicit reporting rather than imputed evidence.

\subsection{Outcomes and denominators to freeze before the mission study}
Let $I_e$ be the set of available native source observations in event $e$, and
$R$ the indices recovered from decoded raw records. Positive-duration event
intervals are half-open; a point event uses its exact native instant. Raw sample
coverage and pre-event context coverage are
\begin{equation}
 C_e = \frac{|I_e\cap R|}{|I_e|},\qquad
 Q_e(w)=\frac{|J_e(w)\cap R|}{|J_e(w)|},
\end{equation}
where $J_e(w)$ contains available source observations during the $w$ time units
before onset. Empty source denominators are explicitly unavailable. If source
observations exist but retention loses every one, coverage is zero rather than
missing. Raw onset error is $\min_{i\in I_e\cap R}(t_i-t_e)$; a missing retained
onset is reported as a failure, not removed from the comparison.

For receiver event recall, the denominator is the source-eligible event set,
shared across policies. A retention policy cannot improve recall by declaring
lost events unscorable. Receiver false alarms and scored-time coverage must be
reported together. The planned primary comparison freezes one receiver model and
threshold across retention policies; model-swap sensitivity is a separate analysis.
Receiver matching, missing-input handling, and summary interfaces must be specified
before the principal mission experiment.

\begin{table}[ht]
\centering
\begin{tabular}{ll}
\toprule
Design variable & Proposed pilot values \\
\midrule
Reporting epoch & One hour \\
Application-byte cap & 256, 512, 1,024, 2,048 bytes per active epoch \\
Pre-trigger ring & 0, 4, 8 native samples \\
Background reservation & 0 or 0.2 of raw-record capacity \\
Evaluation setting & Native Mission1 validation observations \\
\bottomrule
\end{tabular}
\caption{Development grid: the native pilot runs all four caps with an eight-sample
ring and 0.2 background reservation. Other ring/reservation values remain planned
ablations. These are not established mission requirements.}
\end{table}

\subsection{Planned mission study}
The initial source is native ESA Mission1 observations, not the existing
forward-filled grid. The current six-channel subset is a development setting;
a previously untouched mission or independently obtained dataset is needed for a
new generalization estimate. The historical Mission1 test split has already been
examined during detector development and must not be described as pristine.

Fit transforms and capture models on eligible training data, choose configurations
on validation, and freeze them before independent evaluation. Sweep common budget
caps and report actual bytes, event evidence, onset/context preservation, receiver
false alarms, delivery delays, and CPU/memory behavior. Use time-stratified raw
sampling, fixed summaries, threshold-only capture, stronger adaptive sampling,
and a full-archive reference. Include pre/post-context and background-reservation
ablations, capture/receiver model swaps, bursty normal operations, outages, and
faults missed by the trigger. Synthetic faults are stress tests, not replacements
for genuine mission annotations. Pair policy comparisons by event and report
uncertainty without treating correlated frames as independent faults.
